\documentclass[11pt]{article}
\usepackage{fontspec}
\usepackage[french,english]{babel}
\usepackage{amsmath}
\usepackage{amssymb}
\usepackage{amsthm}
\usepackage{mathrsfs}
\usepackage{tikz-cd}
\usetikzlibrary{babel}
\usepackage{enumitem}

\newcommand{\printedpage}[1]{\marginpar{\footnotesize p.~#1}}
\newcommand{\msp}[1]{\marginpar{\footnotesize [#1]}}
\newcommand{\card}{\operatorname{card}}
\newcommand{\redop}{\mathop{\textnormal{red}}\nolimits}
\newcommand{\Hom}{\operatorname{Hom}}
\newcommand{\Ob}{\operatorname{Ob}}
\newcommand{\Ind}{\operatorname{Ind}}
\newcommand{\Ens}{\mathrm{Ens}}
\newcommand{\defeq}{\overset{\text{def}}{=}}
\newcommand{\into}{\lhook\joinrel\longrightarrow}

\setlength{\parindent}{0pt}

\begin{document}

\printedpage{179}
% head: A. Grothendieck : Dérivateurs, Ch. XVIII, Édition des Universités de Montpellier II et Paris VII
% foot: Version 17 décembre 2025

Let us point out at once the

\textbf{Corollary 4.15.10.} \textit{The preceding equivalent conditions are satisfied, if $\pi > \card\mathcal{C}$.}

Indeed, let $\alpha_1$ be the smallest ordinal such that $\card\alpha_1 > \card\mathcal{C}$. Thus $\pi_1 \defeq \card\alpha_1$ is the successor of $\pi_0 \defeq \card\mathcal{C}$, and we have $\pi \geqslant \pi_1$, hence (if $\alpha_0$ is as specified in b$'$ above) $\card\alpha_0 > \pi_1$, hence $\alpha_1 < \alpha_0$. Let (with the notation of b$'$), for $\alpha \in I$ \texttt{[rather} $\alpha \in I_{\alpha_0}$\texttt{]}, i.e.\ $\alpha < \alpha_0$,
\[
M_\alpha \;=\; \varinjlim_{i<\alpha} \underbrace{M_i}_{=\,M}\ ,
\]
we thus see, since the $\alpha_{ji}$ are injective and not surjective, that
\[
\card M_\alpha \;\geqslant\; \card\alpha\ ,
\]
\msp{463}
and that for $\alpha < \beta$, we have
\[
M_\alpha \;\into\; M_\beta \qquad \text{injective}\ .
\]
Let us apply this to the case where $\alpha = \alpha_1$, thus we find
\[
\card M_{\alpha_1} \;\geqslant\; \card\alpha_1 \;=\; \pi_1\ ,
\]
on the other hand, taking $\beta = \alpha_1 + 1$, whence $M_\beta \simeq M$, we find an injection
\[
M_{\alpha_1} \;\into\; M\ ,
\]
whence
\[
\pi_1 \;=\; \card M_{\alpha_1} \;\leqslant\; \card M \;=\; \pi_0\ ,
\]
a contradiction, which establishes b$'$.

\textbf{Corollary 4.15.11.} \textit{Let $\mathcal{C}$ be an essentially small category, $\pi$ a cardinal such that $\pi > \card\redop\mathcal{C}$. Let $\mathcal{M} \defeq \Ind_\pi\mathcal{C}$, then $\mathcal{M} = \mathcal{M}_\pi$, and it is a Karoubi envelope of $\mathcal{C}$.}

We may assume $\mathcal{C}$ reduced, hence $\card\mathcal{C} < \pi$, and we apply the preceding corollary.

Thus, the disputed result 4.8.13 (p.~212) is finally proved in a slightly weaker form, with $\pi > \card\redop\mathcal{C}$ instead of $\pi \geqslant \card\redop\mathcal{C}$. By the same token, the applications of 4.8.13 we have made so far (in particular lemma 4.15.3).

\msp{464}

\subsection*{4.16\quad Examples and questions}

\textbf{4.16.0.} I had thought that every ``reasonable'' functor (or more precisely, every ``logically accessible'' functor) between accessible categories ought to be accessible. This is unfortunately not true, since a composite of two Hom functors
\[
\Ens \longrightarrow \Ens\ , \qquad\quad E \longmapsto \underbrace{\Hom(\Hom(E,I),J)}_{J^{(I^E)}}
\]
\printedpage{180}
% head: A. Grothendieck : Dérivateurs, Ch. XVIII, Édition des Universités de Montpellier II et Paris VII
% foot: Version 17 décembre 2025
is not accessible, if $I$, $J$ are both of cardinality $\geqslant 2$. Taking e.g.\ $I = J = \{0,1\}$, we find the functor
\[
E \longmapsto \mathfrak{P}(\mathfrak{P}(E))\ ,
\]
which is not accessible.

\textbf{4.16.1.} Second disillusionment: I had thought that for every accessible category $\mathcal{M}$, every % [sic] toute catégories -> toute catégorie, singular required after "toute"
strictly full subcategory $\mathcal{N}$ stable under direct factors and ``reasonable'' (e.g. ``logically accessible'') ought to be accessible in $\mathcal{M}$. This is not so, even in the particular case of the inverse image of the accessible subcategory in Ens
\msp{465}
\[
\underbrace{\Ens^*}_{=\,\{E\,\in\Ens\,\mid\, E\neq\emptyset\}} \quad\subset\ \Ens
\]
under an accessible \textit{representable} contrafunctor
\[
\begin{aligned}
\mathcal{M} \quad &\longrightarrow \quad (\Ens)^\circ \\
x \quad &\longmapsto \quad \Hom_{\mathcal{M}}(x,z)\ ,
\end{aligned}
\]
i.e.\ the open (hence strictly full) subcategory $\mathcal{M}_z$ of $\mathcal{M}$ consisting of the objects $x$ such that $\Hom_{\mathcal{M}}(x,z) \neq \emptyset$, is not always accessible in $\mathcal{M}$. Here, I have not arrived at % [sic] n'est par toujours -> n'est pas toujours, misprint "par" for "pas"
as elementary an example (cf.\ appendix) -- I took $\mathcal{M} = (k\text{-comm.\ alg.})$, $k$ an infinite field and $z$ the initial object $k$ of $\mathcal{M}$ -- I found that $\mathcal{M}_z$ is not $\mathrm{L}_\pi$-stable in $\mathcal{M}$, whatever the cardinal $\pi$ % [sic] quelque soit -> quel que soit, standard spelling of the locution
-- but this example used that $\mathfrak{U} = \mathfrak{U}_1$, the smallest universe containing $\mathbf{N}$. I do not know whether for every other universe, one can find a field $k$ of cardinality $\pi \in \mathfrak{U}$ large enough, so that every $\pi$-ultrafilter on a set $E \in \mathfrak{U}$ is trivial. Perhaps
\msp{466}
such an example, where not the construction of the pair $(\mathcal{M},\mathcal{N})$, but the proof of the fact that $\mathcal{N}$ is not accessible in $\mathcal{M}$, appeals to particular hypotheses on the ambient universe, should not be considered as (altogether!) ``logically accessible''. To be really convincing, one would need an example of a more elementary nature.

% margin: a vertical bar is printed in the left margin beside the three lines from "L_pi-stable dans M pour pi assez grand" to "logiquement accessible) serait accessible dans M"
There remains the possibility that there is a metamathematical theorem ensuring that every logically accessible pair $(\mathcal{M},\mathcal{N})$, with $\mathcal{M}$ accessible, and $\mathcal{N}$ a strictly full subcategory $\mathrm{L}_\pi$-stable in $\mathcal{M}$ for $\pi$ large enough, is accessible in $\mathcal{M}$ (\footnote[188]{This is false, even for a closed subcategory, cf.\ 4.16.7.10 below.}). This would apply in particular to the case of a \textit{closed} subcategory: every closed subcategory (that is logically accessible) would be accessible in $\mathcal{M}$. The simplest test example would be that of the closed complement $\mathcal{M}'_z$ of the open subcategory $\mathcal{M}_z$ above (\footnote[189]{\textbf{N.B.} it is accessible in $\mathcal{M}$ if and only if $\mathcal{M}_z$ is! % [sic] elles est -> elle est, agreement with singular verb
}), i.e.\ the category of the $x \in \Ob\mathcal{M}$ such that $\Hom_{\mathcal{M}}(x,z) = \emptyset$, or again
\msp{467}
the inverse image of the accessible closed subcategory
\[
e \;=\; \{\emptyset\} \;\subset\; (\Ens)
\]
reduced to $\{\emptyset\}$, under the representable functor $x \longmapsto \Hom(x,z) : \mathcal{M} \longrightarrow \Ens$. It could thus be that such a category is always accessible in $\mathcal{M}$. But I have been unable to
\printedpage{181}
% head: A. Grothendieck : Dérivateurs, Ch. XVIII, Édition des Universités de Montpellier II et Paris VII
% foot: Version 17 décembre 2025
prove it, even in the particular case above, where $\mathcal{M} = (k\text{-comm.\ alg.})$, $k$ being an infinite commutative field. The difficulty is (in all cases) the following: to find, for fixed $z$, a cardinal $\pi > \pi_0$ such that for every $x \in \Ob\mathcal{M}$, and every system (or \textit{one} system) $\pi_0$-$\pi$-adapted $(x_i)_{i\in I}$ with direct limit $x_0$ ($\pi_0$ a given cardinal, useful for $\mathcal{M}$), \textit{if} $\Hom_{\mathcal{M}}(x,z)$ ($\simeq \varprojlim \Hom_{\mathcal{M}}(x_i,z)$) is empty, then $\exists\ i_0 \in I$ such that $\Hom_{\mathcal{M}}(x_{i_0},z)$ is already empty. In the example $\mathcal{M} = (k\text{-comm.\ alg.})$, one must therefore prove that one can find a cardinal $\pi$
such that for every $k$-algebra $A$ with $\Hom_{\mathcal{M}}(A,k) = \emptyset$, there exists a sub-$k$-algebra $A_i$ of % [sic] pour tout k-algèbre -> pour toute k-algèbre, "algèbre" is feminine
cardinality $\leqslant \pi$, such that $\Hom_k(A_i,k) = \emptyset$ (\footnote[190]{This is false, cf.\ marginal annotation p.~466.}). I tend to believe that this must be true (and even without
\msp{468}
assuming that $k$ is a field, and also taking the condition $\Hom_{\mathcal{M}}(A,Z) = \emptyset$, where $Z$ is an arbitrary object of $\mathcal{M}$). This sort of question seems surprisingly delicate \dots % [sic] étonnement -> étonnamment, adverb required

\textbf{4.16.2.} It is possible, more generally, that the following question has an affirmative answer: Let
\[
f \;:\; \mathcal{M} \longrightarrow \mathcal{N}^\circ
\]
be an accessible functor, with $\mathcal{M}$, $\mathcal{N}$ accessible categories, and let $\underline{T}$ be a subcategory of $\mathcal{N}$ accessible in $\mathcal{N}$, and such that there exists a cardinal $\pi$ having the following properties: $\mathcal{N}^\circ$ satisfies $\mathrm{L}_\pi$ (i.e.\ $\mathcal{N}$ stable under $\varprojlim$ large relative to $\pi_1$). Then is $\underline{S} \defeq f^{-1}(\underline{T})$ accessible in $\mathcal{M}$ (\footnote[191]{False, with $\mathcal{N} = \Ens$ and a representable functor $\mathcal{M} \longrightarrow \Ens^\circ$, taking $I = \{\emptyset\}$, cf.\ marginal annotations p.~466, 467.})? This generalizes the case $\mathcal{N} = (\Ens)$, $\underline{T} = \{\emptyset\}$, $f$ a functor represented by $z \in \Ob\mathcal{M}$, considered just now. With the same type of example, but taking $\mathcal{N} = (\Ens)^*$
\msp{469}
instead of $\mathcal{N} = \{\emptyset\}$, we see that the additional stability condition on $\underline{T}$ is not
superfluous. (\textbf{N.B.} They are to ensure that $\underline{S}$ is $\mathrm{L}_\pi$-stable in $\mathcal{M}$ for large $\pi$.) % [sic] Elle sont -> Elles sont, agreement with plural verb

\textbf{4.16.3.} Related questions: I had hoped that for an accessible functor between accessible categories % [sic] Questions apparentée -> Question apparentée, agreement
\[
f \;:\; \mathcal{N} \longrightarrow \mathcal{M}\ ,
\]
the $\natural$-essential image, consisting of the $y \in \Ob\mathcal{M}$ which are isomorphic to a direct factor of an object of the form $f(x)$, with $x \in \Ob\mathcal{N}$, would be accessible in $\mathcal{N}$. The example above (not at all elementary!) of the subcategory $\mathcal{M}_z$ of $\mathcal{M}$ (consisting of the $x \in \Ob\mathcal{M}$ such that $\Hom_{\mathcal{M}}(x,z) \neq \emptyset$) shows that this is not so. Indeed, $\mathcal{M}_z$ is the essential image, and also the $\natural$-essential image, of the functor
\[
\mathcal{M}/z \longrightarrow \mathcal{M}\ ,
\]
now $\mathcal{M}/z$ is accessible, and the preceding functor accessible, if $\mathcal{M}$ is accessible (verification by AQT \texttt{[âne qui trotte: at a donkey's trot, i.e. by routine verification]}, cf.\ below, more general statement).

\msp{470}

We get the impression that the reason why we have such counterexamples, is that we have $\natural$-essential images which are not $\mathrm{L}_\pi$-stable in $\mathcal{M}$ for any cardinal $\pi \in \mathfrak{U}$. This leads to modifying the conjecture as follows: let $\pi$ be a cardinal such that $\mathcal{M}$ satisfies $\mathrm{L}_\pi$,
\printedpage{182}
% head: A. Grothendieck : Dérivateurs, Ch. XVIII, Édition des Universités de Montpellier II et Paris VII
% foot: Version 17 décembre 2025
and let $\underline{S}(\pi)$ be the smallest strictly full subcategory of $\mathcal{M}$, $\mathrm{L}_\pi$-stable in $\mathcal{M}$, and containing the essential image of $f$. (One could call $\underline{S}(\pi)$ the $\mathrm{L}_\pi$-essential image.) Then $\underline{S}(\pi)$ is accessible in $\mathcal{M}$. Moreover, the (decreasing) family of the $\underline{S}(\pi)$, for varying $\pi$, is stationary, i.e.\ there exists a cardinal $\pi$ such that $\pi' \geqslant \pi$ implies $\underline{S}(\pi') = \underline{S}(\pi)$. It is then clear that $\underline{S}(\pi)$ is the smallest subcategory of $\mathcal{M}$ that is accessible in $\mathcal{M}$ and contains the essential image of $f$. The same question arises for a functor $g^\circ : \mathcal{N}^\circ \longrightarrow \mathcal{M}$, where $g : \mathcal{N} \longrightarrow \mathcal{M}^\circ$ is an accessible functor, $\mathcal{N}$ and $\mathcal{M}$ accessible. I conjecture that the answer is again affirmative.

\msp{471}

% margin: a hand-drawn curved stroke is printed in the left margin beside the last three lines of the following paragraph ("pleine M de C^ qui soit ..." to "sans pourtant être accessible.")
Let $\mathcal{M}$ be a category satisfying condition L, and admitting a small strictly generating full subcategory $\mathcal{C}$, consisting of accessible objects. We thus have a fully faithful and accessible functor
\[
\mathcal{M} \;\into\; \mathcal{C}^\wedge\ ,
\]
$\mathcal{M}$ $\mathrm{L}_\pi$-stable in $\mathcal{C}^\wedge$ for suitable $\pi$. If $\overline{\mathcal{M}}$ is the essential image, so that $\mathcal{M} \simeq \overline{\mathcal{M}}$, it is thus a strictly full subcategory $\mathrm{L}_\pi$-stable in $\mathcal{C}^\wedge$. So we expect it to be accessible in $\mathcal{C}^\wedge$ (if the situation described is ``logically accessible''), or what amounts to the same, $\mathcal{M}$ to be accessible. (Compare 4.8.1, p.~187.) But it should not be difficult (\footnote[192]{OK by taking a \textit{closed} $\underline{S}$ not accessible in $\mathcal{M}$ of an accessible category $\mathcal{M}$, which we embed in a $\mathcal{C}^\wedge$. But these $\underline{S}$ do \textit{not} have a small strictly generating family! % [sic] ces S n'a -> ce S n'a, plural determiner with singular verb
}), by making a number $\pi_{\mathfrak{U}}$ of choices (where $\pi_{\mathfrak{U}} = \card\mathfrak{U}$), to make examples all the same (``logically inaccessible'') of a strictly full subcategory $\mathcal{M}$ of $\mathcal{C}^\wedge$ that is $\mathrm{L}_\pi$-stable, but not accessible in $\mathcal{C}^\wedge$. This $\mathcal{M}$ thus admits a small strictly generating subcategory consisting of accessible objects (\footnote[193]{This is not clear, for this the inclusion functor would have to be a left adjoint!}) and satisfies $\mathrm{L}_\pi$, without however being accessible.

\msp{472}

\textbf{4.16.4} (\footnote[194]{move down after 4.16.5.}). Let $\mathcal{M}$ be a category such that $\mathcal{M}$ and $\mathcal{M}^\circ$ are accessible. I conjecture that then $\mathcal{M}$ is essentially small. More generally, suppose $\mathcal{M}$ accessible, and that $\mathcal{M}$ admits a small subcategory which is \textit{strictly co}generating. Then is $\mathcal{M}$ essentially small? More generally, if $\mathcal{M}$ admits a small strictly generating subcategory, and also a small strictly cogenerating subcategory, i.e.\ if $\mathcal{M}$ and $\mathcal{M}^\circ$ are both equivalent to full subcategories of a $\mathcal{C}^\wedge$, with $\mathcal{C}$ small, then is $\mathcal{M}$ essentially small?

\textbf{4.16.5.} Let $I$ be a small category, $I_\bullet$ the right cone on $I$ (obtained by adding a final object $e$), $\mathcal{M}$ an accessible category,
\[
\underline{\Hom}_{\square}(I_\bullet,\mathcal{M}) \;\into\; \underline{\Hom}(I_\bullet,\mathcal{M})
\]
the strictly full subcategory of the (accessible) category $\underline{\Hom}(I_\bullet,\mathcal{M})$, consisting of the functors $f$ which make $f(e)$ the direct limit in $\mathcal{M}$ of the $f(i)$ ($i \in \Ob I$). This subcategory is $\mathrm{L}_\pi$-stable for large $\pi$, so there is reason
\msp{473}
to conjecture that it is accessible in $\underline{\Hom}(I_\bullet,\mathcal{M})$. $\Big($ I recall that we saw in 4.14 that this is OK if the $\varinjlim$ of type $I$ are representable in $\mathcal{M}$. And that an affirmative answer, in the case where
\printedpage{183}
% head: A. Grothendieck : Dérivateurs, Ch. XVIII, Édition des Universités de Montpellier II et Paris VII
% foot: Version 17 décembre 2025
\[
\begin{tikzcd}[row sep=small, column sep=small]
 & 1 \\
I = 0 \arrow[ur] \arrow[dr] & \\
 & 2
\end{tikzcd}
\quad ,\qquad\quad \text{thus} \qquad\quad
\begin{tikzcd}[row sep=small, column sep=small]
 & 1 \arrow[dr] & \\
I = 0 \arrow[ur] \arrow[dr] & & 3 \\
 & 2 \arrow[ur] &
\end{tikzcd}
\quad ,
\]
implies that the subcategory $\underline{\text{Ép}}(\mathcal{M})$ in $\underline{\mathrm{Fl}}(\mathcal{M})$ is accessible in $\underline{\mathrm{Fl}}(\mathcal{M})$.$\Big)$

\textbf{N.B.} If $\mathcal{M}$ admits a small \textit{co}generating subcategory, then the question falls under conjecture 4.16.2 concerning an accessible functor $\mathcal{M} \longrightarrow \mathcal{N}^\circ$, here $\underline{\Hom}(I_\bullet,\mathcal{M}) \longrightarrow \underline{\Hom}(I^\circ \times \mathcal{C},\Ens)^\circ \dots$

\textbf{4.16.6.} There remains the question of the relative accessibility of a category $\mathcal{E}$ over a small category $\mathcal{C}$. I now think, thanks to 4.15.9, and its corollary 4.15.10, that the following statement must ``fall out'' by AQT \texttt{[âne qui trotte: at a donkey's trot, i.e. by routine verification]}:

\textbf{Proposition 4.16.6.1.} \textit{Let $f : \mathcal{E} \longrightarrow \mathcal{C}$ be a functor, with $\mathcal{C}$ a small Karoubian category. We assume that $f$ satisfies condition $H' = H$ (1+2+3) of 4.15.1 (page 415), where H 2, H 3 suffice, for every arrow $\lambda$} other than an identity \textit{of $\mathcal{C}$. We assume moreover that for a category $\mathcal{C}'$ of one of the two types we are going to describe, and a functor $\mathcal{C}' \longrightarrow \mathcal{C}$, $\mathcal{E}' = \mathcal{E} \times_{\mathcal{C}} \mathcal{C}'$ is accessible, and $\mathcal{E}' \longrightarrow \mathcal{C}'$ accessible.}

\msp{474}

$\Big($ \textit{We take for $\mathcal{C}'$ the category}
\[
\mathcal{Q} = \left\{ \quad
\begin{tikzcd}[column sep=2.2em]
\xi_0 = \xi^\natural \arrow[r, shift left, "\sigma"] \arrow[dr, "u_0"'] & \xi \arrow[l, shift left, "\rho"] \arrow[d, "u = u_0\rho"] \arrow[loop right, "p"] \\
 & \eta
\end{tikzcd}
\qquad
\begin{array}{rclll}
\rho\sigma & = & \mathrm{id}_{\xi_0} & & \\[2ex]
\sigma\rho & = & p\ , & \text{\textit{whence}} & \left\{\begin{array}{rcl} p^2 & = & p \\ p\sigma & = & \sigma \\ \rho p & = & \rho \end{array}\right. \\[4ex]
u & = & u_0\rho\ , & \text{\textit{whence}} & \left\{\begin{array}{rcl} up & = & u \\ u\sigma & = & u_0 \end{array}\right.
\end{array}
\right.
\]
\textit{(three objects $\xi$, $\xi_0$, $\eta$ and five nonidentity arrows $\sigma$, $\rho$, $p$, $u_0$, $u$, linked by the relations above), or the category (isomorphic to $Q^\circ$)} % Doubt resolved against PDF leaf 25: Q is ordinary mathematical italic, retained.
\[
\mathcal{Q}' = \left\{ \quad
\begin{tikzcd}
\xi_0 \arrow[r, shift left, "\sigma"] & \xi \arrow[l, shift left, "\rho"] \arrow[loop right, "p"] \\
 & \eta \arrow[ul, "u_0"] \arrow[u, "u"]
\end{tikzcd}
\qquad
\begin{array}{lcl}
\left.\begin{array}{rcl} \rho\sigma & = & \mathrm{id}_{\xi_0} \\ \sigma\rho & = & p \end{array}\right\} & \Longrightarrow & \left\{\begin{array}{rcl} p^2 & = & p \\ p\sigma & = & \sigma \\ \rho p & = & \rho \end{array}\right. \\[4ex]
\begin{array}{rcl} u & = & \sigma u_0 \end{array} & \Longrightarrow & \left\{\begin{array}{rcl} pu & = & u \\ \rho u & = & u_0 \end{array}\right.
\end{array}
\right.
\]
\textit{We restrict ourselves to functors $\mathcal{Q} \longrightarrow \mathcal{C}$, $\mathcal{Q}' \longrightarrow \mathcal{C}$ such that $p$ is transformed into an arrow} other than an identity \textit{of $\mathcal{C}$ (\footnote[195]{I am not sure if this is reasonable -- to be checked with care!}) (thus the condition relative to $\mathcal{Q}$, $\mathcal{Q}'$ becomes empty, if $\mathcal{C}$ has no projectors).} $\Big)$

\textit{Under these conditions, $\mathcal{E}$ is accessible, and $f : \mathcal{E} \longrightarrow \mathcal{C}$ is accessible.}

\end{document}
